\section{Reusable domain templates}
\label{sec:templates}

This section describes how common mathematical concepts fit the same small
interface.  The templates are not claims that every theorem in the domain is
easy; they separate the routine finite verifier from the theorem-specific
producer.

\subsection{Finite algebra and combinatorics}

A finite structure is encoded by a natural number or a bounded tuple.  Local
relations---incidence, adjacency, multiplication tables, color constraints,
permutation conditions---are fixed bounded formulas and hence have reflected
checkers.  The preferred workflow is:

\begin{enumerate}[label=(\arabic*)]
\item use structural pairing for a fixed record and bounded arrays for
  variable-length tables;
\item reflect the well-formedness predicate;
\item run bounded search or a certified trace;
\item state the public theorem at arbitrary input terms;
\item eliminate noncanonical auxiliary tables unless the output really needs
  them.
\end{enumerate}

Counting arguments use finite sums, injections, and the checked pigeonhole and
descent modules.  A host enumeration may suggest a certificate, but the object
proof must verify every field and derive the conclusion from the reflected
acceptance equation.

\subsection{Rational and real analysis}

Rational numbers are represented by signed integer numerators and positive
denominators, with equality and order reduced to integer polynomial
identities.  A finite real observation is an interval or rational
approximation with a stated error.

A regular-name presentation for a real \(x\) consists of a PR approximant
\(a(k)\) and an object proof
\[
  |a(k)-a(\ell)|\le 2^{-k}+2^{-\ell}.
\]
Operations are implemented query-first.  Addition propagates the desired
output precision to the two inputs.  Multiplication also requests magnitude
bounds.  Reciprocal and normalization require an apartness resource
\(|x|\ge2^{-r}\).  The reverse demand pass calculates the needed input
precision from \(k,r\).

For an existence proof that supplies only approximate witnesses, use a
relational presentation at each precision and prove cross-precision stability.
Do not infer a regular name from pointwise diameter alone.  A name requires an
explicit compatible selector or a stability theorem strong enough that the
least finite selector is coherent.

\subsection{Algebraic and multivalued data}

Polynomial roots, eigenvectors, logarithm branches, Puiseux expansions, and
local inverse branches are naturally relational near degeneracies.  The
certificate may contain:

\begin{itemize}
\item a factorization or approximate factorization;
\item rational disks or boxes containing clusters;
\item a root-count certificate on boundaries;
\item Newton-polygon data;
\item a separation gap;
\item a permutation or matching between two refinements.
\end{itemize}

At collisions, export a multiset, cluster tree, spectral projector, invariant
subspace, or symmetric polynomial observation.  Prove coherence modulo
permutation or within a metric diameter.  If a named branch is genuinely
required, add a positive gap and a deterministic ordering whose stability is
proved from that gap.  The gap is a resource, not a cosmetic side condition.

\subsection{Compactness, optimization, and topology}

Replace abstract compactness by a finite-net interface:

\begin{enumerate}[label=(\alph*)]
\item a PR list of net points at precision \(k\);
\item a reflected coverage relation;
\item a modulus transferring estimates from a point to a nearby net point;
\item bounded search for the least accepted point or least value code.
\end{enumerate}

The coverage theorem is a producer.  Once supplied, finite minimization is
routine.  If minimizers are nonunique, export the minimum value exactly and a
near-minimizer relation, or prove that all selected minimizers have close
observations.

For topology, use finite complexes, carriers, chain maps, and explicit prism
homotopies.  A local-to-global construction should descend only the observation
needed by the theorem.  A checked spanning tree can transport local data and
accumulate rational edge errors; constructing a global quotient object is
unnecessary unless it appears in the endpoint.

\subsection{Probability and statistics}

Begin with finite rational laws.  Expectation is a finite weighted sum;
variance and moments are polynomial identities after denominator clearing.
For larger spaces, a presentation must supply finite quantizations, tail
bounds, integration moduli, and any anti-concentration data used by a
discontinuous statistic.

At an atom, an exact median or quantile may be nonunique or discontinuous.
Export the interval of valid quantiles or a robust bracket.  A unique exact
quantile can be named only under a gap or anti-concentration premise that a
finite checker can consume.

\subsection{Diophantine and arithmetic geometry}

Given an explicit height bound \(H\), enumerate every code below an effective
bound \(E(H)\), test the Diophantine relation, and return the complete list.
The verifier proves:

\begin{enumerate}[label=(\roman*)]
\item every listed item is a solution;
\item no list entry is duplicated under the chosen code normalization;
\item every solution of height at most \(H\) occurs in the list.
\end{enumerate}

The remaining theorem is exactly
\(P(x)\to h(x)\le H\).  It may be a Runge estimate, a descent theorem, a Baker
bound, a supplied covering list, or an open effective producer.  Mordell--Weil
and Selmer computations likewise separate finite class arithmetic from the
existence of rational representatives on coverings.  The class verifier can
be complete while the representative producer remains a residual.

\subsection{Categories and abstract structures}

For a fixed finite presentation, categorical constructions are finite tables
with equation checkers.  Composition, finite limits and colimits, finite Kan
extensions, and finite descent can therefore be treated by bounded reflection.
The key restriction is ``fixed finite presentation.''  A universal evaluator
for arbitrary categories or arbitrary PR programs is neither needed nor
available.  Functoriality and naturality are public object derivations over the
fixed codes.
